\documentclass[11pt,a4paper]{article}
\usepackage[margin=25mm]{geometry}
\usepackage{amsmath,amssymb,amsthm,booktabs,longtable,hyperref}
\newtheorem{theorem}{Theorem}
\newtheorem{lemma}[theorem]{Lemma}
\newtheorem{proposition}[theorem]{Proposition}
\newtheorem{corollary}[theorem]{Corollary}
\newcommand{\Sh}{\operatorname{Sh}}
\newcommand{\Sym}{\operatorname{Sym}}
\newcommand{\Aut}{\operatorname{Aut}}
\title{Witnesses and orbits in finite incidence geometries}
\author{Private common-theory assembly; independent review pending}
\date{}
\begin{document}\maketitle
This draft assembles complete previously reviewed proof components into
one common theory. A geometric shadow intersection must retain a flag
witness; a subgroup intersection identity for distinguished generators
is a different property. The elementary covering formulation, finite
coset normalization, all-residue tests and sparse realization below
make this distinction explicit. The common forest and thin-geometry
arguments then support the two uniform applications reported in L096.
No new theorem priority, whole-source completion or measurement is
claimed by this unfinished assembly.

The source orbit and witness reductions are due to Jacobs and Leemans,
\emph{An algorithmic analysis of the intersection property}, LMS Journal
of Computation and Mathematics 7 (2004), 284--299, Section3. The block,
stabilizer and transporter arguments are standard. The exact migration
ledger distinguishes inherited proofs from the present reorganization;
all paid common and retained classification premises remain explicit.
\section{Flag-compatible covering}
Let $X$ be a finite set partitioned into type classes $X_i$, with a
reflexive symmetric incidence relation $\sim$ such that distinct objects
of the same type are not incident. A flag is a pairwise incident set,
including the empty set. Define the common shadow of any set $U$ by
\[
 \sigma_i(U)=\{y\in X_i:y\sim u\text{ for every }u\in U\}.
\]
The query set need not itself be a flag. The geometric intersection test
uses $U=\{x\}\cup F$, where $F$ is a flag. It asks whether the nonempty
set $S=\sigma_i(U)$ is the shadow of a flag incident with every object
of $U$. A singleton $S$ has its unique object as a witness; the whole
set $X_i$ has the empty flag as a witness.

\begin{theorem}\label{cover}
Suppose $|S|\geq2$, and put
\[
 E=\{v\in X\setminus X_i:v\sim u\ (u\in U),\ v\sim s\ (s\in S)\}.
\]
A witness flag exists if and only if some pairwise incident set $C\subseteq E$
meets every exclusion set
$\{v\in E:v\not\sim y\}$, $y\in X_i\setminus S$.
\end{theorem}
\begin{proof}
A witness cannot contain an object of type $i$, since that would reduce
its shadow to at most one object. Its incidence with $U$ and its shadow
$S$ therefore place every member in $E$. Each $y$ outside its shadow
fails incidence with at least one member, giving the exclusion condition.
Conversely, every eligible member is incident with $S$, so an eligible
flag has shadow containing $S$. The exclusion condition removes every
object outside $S$, proving equality. Eligibility also makes the flag
incident with $U$. The argument includes $C=\varnothing$ when there are
no objects to exclude.
\end{proof}
Equivalently introduce a Boolean variable $z_v$ for each $v\in E$ and
require
\[
 z_v+z_w\leq1\quad(v\ne w,\ v\not\sim w),\qquad
 \sum_{\substack{v\in E\\v\not\sim y}}z_v\geq1
       \quad(y\in X_i\setminus S).
\]
The first constraints include incompatible objects of the same type.
Thus the criterion is an exact finite covering problem with flag
compatibility retained. A solution certifies a witness; a heuristic
failure is not a negative proof. The exhaustive finite algorithms below
preserve the entire witness space.

Dropping compatibility is unsound, even for a rank-three geometry. Take
$A=\{0,1,2,3\}$, $B=\{b_0,b_1\}$ and $D=\{d_0,d_1\}$. Make every
$B$ object incident with every $D$ object, and put
\[
 N_A(b_0)=N_A(d_0)=\{0,1,2\},\qquad
 N_A(b_1)=N_A(d_1)=\{0,1,3\}.
\]
Every vertex and cross-type edge extends to a chamber, so this is a
geometry. For $x=b_0$, $F=\{b_1\}$, the target shadow is $S=\{0,1\}$.
The eligible set is $\{d_0,d_1\}$ and its common shadow is exactly $S$,
but its two objects have the same type. No eligible flag has shadow $S$:
the empty flag has four target objects and each singleton has three.
The covering formulation and this obstruction reuse the project's
previously reviewed proof; they are not a replacement claim to the
original source's orbit algorithm.

\section{Typed cosets and flag transitivity}
The following construction is Tits' coset construction, recalled as
Theorem~3.1 of Jacobs--Leemans. No finiteness of the group is needed here.
Let $I$ be a finite type set, $G$ a group, and $(H_i)_{i\in I}$ subgroups.
Take the \emph{disjoint typed union} of the right-coset sets
$H_i\backslash G$; the type remains part of an object even when two
subgroups happen to be equal. Declare two typed cosets incident when
their underlying subsets of $G$ intersect.

\begin{proposition}\label{tits}
This is a reflexive symmetric incidence structure, with no incident
distinct objects of the same type. The cosets $(H_i)_{i\in I}$ form a
chamber. Right multiplication by $G$ preserves incidence and types and
is transitive on flags of each prescribed type of size at most two.
It need not be transitive on larger flags, nor need every flag extend
to a chamber without an additional hypothesis.
\end{proposition}
\begin{proof}
Right cosets of a fixed subgroup are either equal or disjoint. This
gives the same-type assertion; reflexivity and symmetry follow from
intersection. The displayed cosets all contain the identity and include
one object of every type. Multiplication on the right is a bijection
of $G$ preserving intersections. Every object is a translate of its
standard coset. If $H_i g$ and $H_j h$ are incident, choose a common
element $a$; the two cosets equal $H_i a,H_j a$, respectively. Thus this
pair is the translate by $a$ of the standard pair. The empty flag is
fixed. These arguments prove exactly the asserted transitivity.
\end{proof}

In the finite algorithms below, full flag transitivity is proved by
checking all eligible-extension equalities, not inferred from this
pair-transitivity statement. A nontrivial kernel of the coset action
is allowed. Inversion identifies this right-coset convention with the
left-coset convention used in the forest application, so no extra
geometric premise is introduced when changing conventions.

\section{Orbit and witness recognition}
An incidence geometry has a finite set of types $I$, reflexive symmetric
incidence, and no incident distinct objects of the same type. Every flag
extends to a chamber. Suppose a finite group $G$ acts transitively on flags
of each type. Fix a chamber $C=\{p_i:i\in I\}$, use right actions, and write
\[
 H_i=G_{p_i},\quad H_T=\bigcap_{t\in T}H_t,\quad H_\varnothing=G,
 \qquad C_T=\{p_t:t\in T\}.
\]
The $i$-shadow $\Sh_i(F)$ is the set of type-$i$ objects incident with every
member of $F$. This geometric shadow property is distinct from the
subgroup intersection property for distinguished generators of a C-group.

\begin{lemma}\label{orbit}
For every $i,T,h$ we have
\[
             \Sh_i(C_T^h)=(p_i^{H_T})^h.
\]
\end{lemma}
\begin{proof}
If $i\in T$, both sides are the singleton $\{p_i^h\}$.
Otherwise a type-$i$ object $y$ incident with $C_T$ completes it to a flag
of type $T\cup\{i\}$. Flag transitivity gives an element carrying
$C_{T\cup\{i\}}$ to that flag while fixing each member of $C_T$.
It belongs to $H_T$ and carries $p_i$ to $y$. Conversely every element of
$H_T$ fixes $C_T$, so its image of $p_i$ is incident with $C_T$.
Apply $h$ to obtain the stated equality.
\end{proof}

The objects of type $i$ are identified with right cosets $H_i\backslash G$.
Thus the lemma uses one permutation action for each type. All shadows in
that action are orbits of intersections of the chamber stabilizers.
The recognition theorem in the current common theory for finite geometries
can establish the required flag transitivity: for every $i\notin T$, the
intersection $\bigcap_{t\in T}p_i^{H_t}$ must equal $p_i^{H_T}$.
Induction on $|T|$ proves that this condition is sufficient, starting with
the empty flag. It checks the entire set of eligible extensions, not a
sample of chambers.

\begin{theorem}\label{shadowtest}
Fix a type $i$. The following condition is equivalent to the geometric
shadow intersection condition at type $i$: for every element $x$ and flag
$F$, the intersection $\Sh_i(x)\cap\Sh_i(F)$ is empty or equals
$\Sh_i(F')$ for a flag $F'$ incident with both $x$ and $F$.

For each $j\ne i$, put $K=H_i\cap H_j$. For every nonempty
$T\subseteq I\setminus\{i\}$, choose representatives $h$ of
\[
                  (H_i\cap H_T)\backslash H_i/K.
\]
For each such representative form
\[
                  X=p_i^{H_j}\cap(p_i^{H_T})^h.
\]
If $|X|>1$, there must be a subset $U\subseteq I\setminus\{i\}$
and an element $q\in K$ such that
\[
 (p_i^{H_U})^q=X,
 \qquad p_u^q\in(p_u^{H_T})^h\quad\hbox{for every }u\in U.
\]
For a fixed $U$ it suffices to take one representative of each right
coset of $K\cap H_U$ in $K$.
\end{theorem}
\begin{proof}
Cases where $x$ has type $i$, $F$ contains type $i$, or $F$ is empty have
an immediate witness whenever the intersection is nonempty. Otherwise
choose $p$ in the nonempty intersection. The pair $\{p,x\}$ is a flag of
types $\{i,j\}$; normalize it to $\{p_i,p_j\}$. The normalized $F$ lies
in the residue of $p_i$. Flag transitivity on $T\cup\{i\}$ expresses it
as $C_T^h$ with $h\in H_i$. Its stabilizer inside $H_i$ is $H_i\cap H_T$.
The remaining choices of normalization act by $K$, giving precisely the
displayed double cosets. Lemma~\ref{orbit} identifies its intersection
of shadows with $X$.

The singleton case has the witness consisting of its unique type-$i$
object. If $|X|>1$, a witness cannot contain type $i$. Let its type be
$U$. Because $p_i\in X$, it is incident with $p_i$, and by assumption it
is incident with $p_j$. Hence its union with $\{p_i,p_j\}$ is a flag.
Flag transitivity expresses the witness as $C_U^q$ for $q\in K$.
Two elements of $K$ give the same such flag exactly when they are in
the same right coset of $K\cap H_U$. Lemma~\ref{orbit} gives the first
displayed condition. Incidence with $C_T^h$ says that, for every $u\in U$,
its type-$u$ object belongs to $\Sh_u(C_T^h)$, giving the second.

Conversely $C_U^q$ is a flag. Since $q$ fixes $p_i$ and $p_j$, it is
incident with those two objects. The second displayed condition makes
it incident with $C_T^h$, and the first identifies its $i$-shadow with
$X$. Transporting it back gives the required original witness. Every
normalization and every orbit representative preserves these conditions,
which proves sufficiency as well as necessity.
\end{proof}

There are two useful exact shortcuts. The candidate type $U$ can be
discarded unless $|p_i^{H_U}|=|X|$. Also, if $p_j$ is incident with
$C_T^h$ and $|X|$ is the smaller of the two original shadow sizes, one
of the original flags is already a witness. Both are consequences of
the theorem and preserve its complete search space.

The source also omits certain candidate types in its implementation.
These omissions have a direct justification. If a witness $F'$ is
incident with $F$ and its type $U$ is contained in the type $T$ of $F$,
the same-type incidence axiom forces every member of $F'$ to be the
corresponding member of $F$. Hence $F'\subseteq F$ and
$\Sh_i(F')\supseteq\Sh_i(F)$. Such a type cannot realize an intersection
strictly smaller than $\Sh_i(F)$. If the intersection has the smaller
of the two original sizes but $x$ is not incident with $F$, a candidate
of type $U=T$ would have to equal $F$ and so would fail incidence with
$x$. It can also be omitted. These prove the two type-pruning steps in
the source Figure~1. The present complete search may retain those
impossible candidates; doing so changes time but not the answer.


\section{Every residue and every rank}
For $J\subseteq I$, let $R_J$ be the residue of the base flag $C_J$.
Its types are $I\setminus J$. Empty residues are not tested for an
existential type condition.
\begin{theorem}\label{residue}
Suppose the coset system is flag transitive. The map
\[
 (H_J\cap H_i)h\longmapsto H_i h,\qquad h\in H_J,
 \quad i\in I\setminus J,
\]
is an isomorphism of the coset geometry inside $H_J$ with $R_J$.
In particular $H_J$ is flag transitive on $R_J$.
For each $k\in\{2,\ldots,|I|\}$ it suffices to test precisely the base
residues $R_J$ with $|I\setminus J|=k$. On each of them apply
Theorem~\ref{shadowtest} with group $H_J$ and parabolics $H_J\cap H_i$.
All type tests in all these residues give $(\mathrm{IP})_k$;
at least one type test in each residue gives $(\mathrm{WIP})_k$.
The conjunction over $k=2,\ldots,|I|$ gives respectively IP and WIP.
\end{theorem}
\begin{proof}
An object of type $i$ incident with $C_J$ extends that flag. Flag
transitivity therefore writes it as $H_i h$ for $h\in H_J$. Two such
representatives give the same object exactly when they give the same
right coset of $H_J\cap H_i$ inside $H_J$. This proves object bijectivity.
If the small cosets intersect, their common representative gives incidence
of their images. Conversely, two incident images together with $C_J$
form a flag. Flag transitivity gives a representative in $H_J$ for both
images, so their small cosets intersect. The same argument, applied to
an arbitrary flag of the residue together with $C_J$, proves flag
transitivity under $H_J$.

Every flag of type $J$ is a translate of $C_J$, so all its residues are
isomorphic with their type labels retained. This proves that one base
residue per subset is sufficient; it does not identify different subsets
or assume their residues are isomorphic. Inside a fixed residue the
preceding theorem is an equivalence at each remaining type. Universal
quantification of those types is IP at that rank; the existential type
for WIP is chosen separately in each residue. These are precisely
[JL, Lemmas 2.1--2.2]. Rank-one shadow intersections are empty, singletons,
or the whole sole type, with the corresponding immediate witnesses.
Rank zero has no available type and is excluded, as in those lemmas.
\end{proof}
There are $[G:H_J]$ flags of type $J$, so the report records that many
actual residues represented by each base residue. It computes every
remaining type, even when one failed type would already refute IP.
Thus its WIP output does not silently use a fixed global type.

\section{Geometry hypotheses are separate checks}
The supplementary rows are not accepted because their total matches a
published table. They define coset incidence structures, which need not
be flag transitive. We first check the orbit equality preceding
Theorem~\ref{shadowtest} for all $i,T$. This condition is necessary by
Lemma~\ref{orbit}. For sufficiency, normalize an existing flag of type
$T$ inductively. Every object eligible to extend it belongs to the
intersection of the singleton shadows, hence to the $H_T$ orbit by the
checked equality. It follows that each extended flag is a translate of
the corresponding base flag. Induction covers every flag type, and
in particular every flag extends to a chamber.

Write $K=H_J$ and $P_i=K\cap H_i$ for a residue. The following tests
are the actual source hypotheses, with kernels of actions allowed.
\begin{lemma}\label{filters}
For a flag-transitive coset geometry the following conditions are exact.
\begin{enumerate}
\item Firmness holds if and only if $[H_{I\setminus\{i\}}:H_I]\ge2$
for every $i$.
\item Residual connectedness holds if and only if
$\langle P_i:i\in I\setminus J\rangle=K$ for every residue of rank
at least two.
\item Residual weak primitivity holds if and only if in each residue of
positive rank at least one transitive action on $P_i\backslash K$ is
primitive of degree at least two.
\item The local condition $(2T)_1$ holds if and only if, in every
rank-one residue, its point stabilizer has exactly two orbits on its
objects and there are at least two objects.
\end{enumerate}
\end{lemma}
\begin{proof}
The objects completing the base panel missing type $i$ are the cosets
of $H_I$ in $H_{I\setminus\{i\}}$. Flag transitivity covers every panel,
proving the first claim. For the second, a coset path starting at the base
chamber can multiply its representative successively by elements of the
$P_i$. Conversely each such multiplication is realized by a path through
the object fixed by that subgroup. Hence the connected component of the
base chamber is exactly the orbit of their join, and all components are
its cosets. The argument is applied inside each $K$; it is valid for
rank two as well as higher ranks. The third statement is the definition
of weak primitivity applied to each residue. Equivalently $P_i$ must be
maximal proper in $K$: blocks containing the base coset correspond to
intermediate subgroups. This equivalence does not require a faithful
action. Finally a transitive action is two-transitive exactly when the
stabilizer of a point is transitive on the other points. Its two orbits
are then that point and its complement. Again this depends on the
permutation image, not on a faithful ambient action.
\end{proof}
These are the panel and coset-path arguments already used in the accepted
Hall--Janko census [D]. The implementation tests the entire action, not
only subgroup orders or a presumed diagram.

\section{Rank two without a diagram assumption}
In a rank-two incidence structure the only nonautomatic shadow test
compares the neighborhoods of two distinct objects on the same side.
At a fixed shadow type $i$, it passes exactly when every such intersection
has size zero, one, or the whole type-$i$ set. Indeed a witness of type
$i$ has singleton shadow; a witness of the other type cannot be incident
with both distinct objects of that type. The empty flag has the whole
shadow set. These are all possible witness flags. In particular comparing
an object with itself is automatic and is excluded from this pair test.

Both type tests hold if and only if the graph has no four-cycle or is
complete bipartite. If a four-cycle is present, its two vertices on one
side have at least two common neighbors, hence all vertices of the other
side. Every pair on that other side consequently has at least two common
neighbors and must meet the entire first side. This proves completeness.
The converse follows immediately. This is the accepted direct proof
[Cover], recovering the partial-linear-space/generalized-digon criterion
used in [JL, Section 3]. We still record the two directional tests
separately, and do not replace the weak condition by an assumed diagram.

For a coset action, fix the first object of the opposite type $j$.
Representatives of $P_j\backslash K/P_j$ cover its possible second
objects. Discard the identity double coset, which represents the same
object. Compute the two neighborhoods by their actual coset-action
orbits. Their intersections then perform the criterion just proved.


\section{A sparse realization of the common shadow theorem}
For a finite permutation group $K$ and subgroups $P_1,\ldots,P_r$, write
$P_T=\bigcap_{t\in T}P_t$ and $P_\varnothing=K$. The $i$-shadow of the
base flag of type $T$ is $\{P_i x:x\in P_T\}$. The original implementation
realizes the whole right-coset action of degree $[K:P_i]$. The following
representation avoids constructing a large universe when only small
shadows are used.
\begin{proposition}
For nonempty $T$, take a right transversal of $P_i\cap P_T$ in $P_T$ and
replace each representative $x$ by a canonical representative of $P_i x$.
This set represents the entire $i$-shadow without duplicates. Its right
translate by $g$ is obtained by replacing each representative $x$ by the
canonical representative of $P_i xg$. The empty flag may instead be
represented by the cardinality $[K:P_i]$, provided its universal-set role
is handled explicitly as below.
\end{proposition}
\begin{proof}
For $x,y\in P_T$, equality $P_i x=P_i y$ holds exactly when
$xy^{-1}\in P_i\cap P_T$. Thus the transversal gives exactly one
representative per shadow object. Canonicalization changes only its name.
If $x'=hx$ with $h\in P_i$, then $P_i x'g=P_i xg$, so the translated
representative is well defined and realizes the same right action.
Intersections, containment, and membership are consequently unchanged.

The empty flag has the full object set as shadow. Its flag-extension
identity is tautological and is counted as such. In a witness test, the
computed intersection is already a subset of the $i$-object set. It is
therefore the empty flag's shadow exactly when its cardinality is
$[K:P_i]$; the empty flag is automatically incident with both tested flags.
When this equality fails, the empty flag cannot witness the intersection,
and every nonempty witness is tested by the displayed representation.
No intersection with a fictitious empty set is performed. The same degree
is used in the rank-two complete-bipartite exception.
\end{proof}
The normalization, double-coset enumeration and flag-constrained witness
search are exactly those of the accepted common theorem. Only the
representation of each finite shadow changes. Trivial panel stabilizers
are allowed: the subgroup generated by an empty generator list is formed
with the identity included, so a two-object panel is still recognized.
The original accepted implementation remains unchanged.


\section{Primitive actions and exact family transport}
\subsection{Primitivity without constructing a large coset action}
Let $H\leq G$ be finite. The transitive action of $G$ on its right cosets
of $H$ is primitive of degree at least two if and only if $H<G$ and
\[
 \langle H,x\rangle=G
 \quad\hbox{for every representative }x\notin H
 \hbox{ of an }(H,H)\hbox{-double coset in }G.                 \tag{*}
\]
This is the standard block--overgroup correspondence together with an
explicit double-coset test. Here is the complete argument. A block
containing the base coset $H$ is the orbit of $H$ under a subgroup
$K$ with $H\leq K\leq G$: its setwise stabilizer is $K$, and transitivity
shows the block is exactly $\{Hk:k\in K\}$. Conversely this set is a
block for each such $K$. Hence degree-at-least-two primitivity is exactly
maximal properness of $H$. If $g=h_1xh_2$, with $h_1,h_2\in H$, then
$\langle H,g\rangle=\langle H,x\rangle$. A proper intermediate subgroup
therefore contains an element outside $H$ and supplies a representative
violating (*); conversely any violating representative gives the proper
intermediate subgroup $\langle H,x\rangle$. The argument does not require
a faithful coset action or a core-free stabilizer. The case $H=G$ must be
rejected separately: its degree-one action is outside the convention and
would otherwise satisfy (*) vacuously.

The finite implementation enumerates the complete double-coset roster,
forms each indicated generated subgroup and compares its order with $|G|$.
Every negative nontrivial case records the outside element and actual
intermediate subgroup generators. For a residue it tests types in
increasing coset degree until a primitive type is found, or exhausts all
types. This preserves the existential weak-primitivity quantifier.
Applying it to all base residues is legitimate by the separately checked
flag transitivity; the residue common theorem supplies their actual
stabilizers and translates. No previously classified geometry predicate
is used by this eligibility test.

\subsection{Exact family transport with necessary order filters}
Let $A\leq\operatorname{Aut}(G)$ be a finite, explicitly computed
transformation group, and let $\mathcal H=(H_1,\ldots,H_n)$ and
$\mathcal K=(K_1,\ldots,K_n)$ be ordered families. We seek an element of
$A$ carrying the first family to a permutation of the second. Put
$H_J=\bigcap_{j\in J}H_j$, with $H_\varnothing=G$, and define $K_J$
similarly. If a proposed type permutation sends $H_i$ to $K_{\pi(i)}$,
then necessarily
\[
 |H_J|=|K_{\pi(J)}|\qquad(J\subseteq\{1,\ldots,n\}).       \tag{**}
\]
Indeed an automorphism commutes with intersection and preserves order.
Thus a failed equality rejects that type permutation, and a failed
comparison of the two multisets of intersection orders rejects every
permutation. Satisfying (**) never certifies equivalence.

Fix a surviving target tuple $\mathcal T$ and any order
$i_1,\ldots,i_n$ in which to place its entries. Start with $C=A$ and
$\mathcal H' =\mathcal H$. At step $k$ find an actual transporter
$x\in C$ with $x(H'_{i_k})=T_{i_k}$. If none exists, reject this target.
Otherwise replace every $H'_i$ by $x(H'_i)$ and replace $C$ by the
setwise stabilizer of $T_{i_k}$ in $C$. After the last step no further
stabilizer is needed. Success supplies a product of the computed
automorphisms taking the entire family to the target.

For completeness, after the first $k-1$ placements all possible further
solutions differ from the current partial transporter by an element of
the pointwise stabilizer of the already placed \emph{subgroups}. If $x$
is one solution to the next placement, all such solutions differ from
$x$ by an element fixing its target subgroup. Hence passing to the stated
stabilizer loses no solution. Induction proves both soundness and
completeness, for every placement order. The stabilizer after $k$
placements depends only on $A$ and the first $k$ target subgroups, not on
the input family or the transporters chosen. It can therefore be cached
for repeated comparisons against the same target and placement order.
In this implementation larger target subgroups are placed first, with
indices breaking ties. This changes computation time, not the criterion.
Literal equality of complete tuples is an immediate successful case.

\subsection{Reducing automorphisms to simultaneous conjugacy}
There is a useful version when a complete set $B$ of representatives
for $\operatorname{Aut}(G)/\operatorname{Inn}(G)$ is available. For each
$\beta\in B$ and each target type permutation, apply $\beta$ to the
whole input family. The remaining question is simultaneous conjugacy
by one element of $G$. In the preceding procedure start with $C=G$;
replace subgroup transport under automorphisms by conjugation, and
replace each target stabilizer by its normalizer in $C$.

Every automorphism is an inner automorphism composed with one tested
representative, since the inner automorphisms form a normal subgroup.
For a fixed representative, the same transporter-coset argument proves
that the successive normalizers preserve exactly all possible remaining
conjugators. A positive result records the representative's images of
the ambient generators, one actual conjugating element, and the type
permutation; together these identify the entire family. Neither a
splitting of the automorphism group nor a representative of the same
order as its quotient image is required. In particular, when the inner
subgroup has index two, any verified noninner automorphism together
with the identity is a sufficient complete transversal. This is a
standard quotient and transporter reduction; its implementation below
is not a new classification theorem.


\section{Forest subgroups and their incidence geometry}
Let $T$ be a finite tree on a set $\Omega$ of $N\geq2$ vertices, with edge
set $E$. For $F\subseteq E$, let $K_F$ be the subgroup of $\Sym(\Omega)$
generated by the endpoint transpositions of the edges in $F$. Set
$H_e=K_{E\setminus\{e\}}$. The elements of type $e$ in $\Gamma(T)$ are
left cosets $gH_e$. Two such elements are incident when their cosets
intersect; incidence is reflexive, and distinct elements of one type are
not incident. A flag is a set of pairwise incident elements. The residue
of a flag consists of its incident elements of the remaining types.

\begin{lemma}\label{forest}
If $C$ ranges over the components of $(\Omega,F)$, then
\[
 K_F=\prod_C\Sym(C),\qquad K_F\cap K_D=K_{F\cap D}.
\]
The coset system $\Gamma(T)$ is a thin, residually connected geometry on
which $\Sym(\Omega)$ is flag-transitive. The residue at the standard flag
of types $E\setminus F$ is the join of the geometries of the nontrivial
components of $(\Omega,F)$.
\end{lemma}
\begin{proof}
Transpositions along a connected graph generate its full symmetric group:
conjugating successive edge transpositions along a path produces the
transposition of its endpoints. Distinct components have disjoint
supports. A permutation in both $K_F$ and $K_D$ preserves every block of
the common refinement of their component partitions. Two vertices are
in such a block exactly when their unique tree path belongs to both
forests, hence to $F\cap D$. This proves the two group formulas.

Root $T$, and let $A_e$ be the component of $T-e$ away from the root.
The coset $gH_e$ identifies with the subset $gA_e$ of $\Omega$.
For distinct types, incidence is precisely
\[
 |X_e\cap X_f|=|A_e\cap A_f|.
\]
Indeed these intersection sizes and the two fixed subset sizes determine
all four atoms of the pair; equal atom sizes give a permutation carrying
one pair to the other. The $A_e$ form a laminar family. Thus pairwise
incidence forces exactly their nesting and disjointness relations on the
$X_e$. Adjoin $\Omega$ above the selected sets and subtract each
container's maximal selected proper subcontainers. The resulting atoms
partition $\Omega$ and have the same sizes on the standard and target
sides. Bijections of corresponding atoms combine to one permutation
realizing the entire flag. Every flag therefore extends to a chamber,
and flags of each fixed type form one orbit.

The stabilizer of the standard flag of types $E\setminus F$ is $K_F$.
Its maximal parabolics for the remaining types are $K_{F\setminus\{e\}}$.
Their direct product descriptions identify the residue with the claimed
join; isolated vertices contribute trivial factors. The chamber
stabilizer is $K_\varnothing=1$, and each panel stabilizer has order two,
so the geometry is thin. In any residue of rank at least two, the
maximal parabolics collectively contain every remaining edge
transposition, and hence generate the residue group. Their coset
incidence graph is connected: products of subgroup elements give paths
between standard and arbitrary cosets. This proves residual connectedness.
\end{proof}

For an element type $i$, write $\sigma_i(F)$ for the set of elements of
type $i$ incident with every member of $F$. The \emph{intersection
property} requires that for every element $x$ and flag $F$, a nonempty
$\sigma_i(x)\cap\sigma_i(F)$ be $\sigma_i(F')$ for a flag $F'$ incident
with $x$ and $F$. This is required in the geometry and in every residue.
It is distinct from the subgroup intersection identity of Lemma~\ref{forest}.

\section{Explicit witnesses for Boolean shadows}
Let $\mathcal B(\Omega)$ have the nonempty proper subsets of $\Omega$ as
elements, cardinality as type, and comparability by inclusion as incidence.

\begin{lemma}\label{boolean}
Fix $0<k<N$. For any family $\mathcal U$ of elements of
$\mathcal B(\Omega)$, a nonempty common $k$-shadow has a witness flag of
at most two elements, incident with every member of $\mathcal U$.
\end{lemma}
\begin{proof}
If $\mathcal U$ contains a $k$-subset $X$, every common shadow element
must equal $X$. Nonemptiness says that $X$ is incident with the whole
family, so $\{X\}$ is the required flag. Otherwise put
\[
 L=\bigcup_{\substack{X\in\mathcal U\\|X|<k}}X,
 \qquad V=\bigcap_{\substack{X\in\mathcal U\\|X|>k}}X,
\]
with empty union $\varnothing$ and empty intersection $\Omega$.
The common shadow is exactly
\begin{equation}\label{interval}
 \{Y\subseteq\Omega: |Y|=k,\ L\subseteq Y\subseteq V\}.
\end{equation}
Nonemptiness gives $L\subseteq V$ and $|L|\leq k\leq|V|$.
If either cardinality bound is equality, the shadow has a unique element,
which again supplies a singleton flag incident with the whole family.
In the remaining case use $\{L,V\}$, omitting $\varnothing$ and
$\Omega$. Its members are proper subsets of different sizes and form
a flag. Every smaller query member is contained in $L$, and every
larger one contains $V$, so both witness elements are comparable with
every query member. Its shadow is precisely~\eqref{interval}, including
when one or both boundary sets were omitted.
\end{proof}

\begin{lemma}\label{boolean-residues}
The geometry $\mathcal B(\Omega)$ satisfies the intersection property
in all residues. Joins of geometries with that property also satisfy it.
\end{lemma}
\begin{proof}
Apply Lemma~\ref{boolean} to $\mathcal U=\{x\}\cup F$ for the property
in $\mathcal B(\Omega)$. A fixed chain, completed by the formal endpoints
$\varnothing$ and $\Omega$, divides its residue into the proper intervals
between successive members. An interval $(A,B)$ is the Boolean geometry
on $B\setminus A$, via $X\mapsto X\setminus A$; the types are shifted
by $|A|$. All elements in different intervals are incident, so the residue
is their join. Empty intervals add no types. A rank-one factor has only
empty, full or singleton shadows; the empty flag witnesses a full
shadow, and its unique point witnesses a singleton.

For completeness, in a join fix a type $i$ in one factor. Only query
elements of that factor constrain its shadow. If $x$ belongs to it, use
the factor's intersection property with the part of $F$ in that factor.
If $x$ belongs to another factor, the part of $F$ in the target factor
is itself a witness, including when it is empty. The witness is
automatically incident with all query elements in other factors.
Residues of a join are joins of factor residues, proving the assertion
also at every residue.
\end{proof}

\begin{theorem}\label{path}
The geometry $\Gamma(T)$ satisfies the intersection property if and only
if $T$ is a path.
\end{theorem}
\begin{proof}
For a path, root at one endpoint. Its cuts $A_e$ form a chain with sizes
$1,\ldots,N-1$. The subset description in Lemma~\ref{forest} identifies
its incidence with comparability, so $\Gamma(T)=\mathcal B(\Omega)$.
Lemma~\ref{boolean-residues} proves the positive assertion.

A tree which is not a path has three edges incident at a common vertex.
Keep precisely those edges in Lemma~\ref{forest}. The resulting residue
is the geometry of the claw, with three types $a,b,c$, each a copy of
$\{1,2,3,4\}$. Elements of different types are incident exactly when
their labels differ. Take $x=(b,1)$ and $F=\{(b,2)\}$, with target
type $a$. Their common shadow is $\{(a,3),(a,4)\}$. A witness cannot
contain an $a$-element, which would give at most a singleton shadow,
or a $b$-element, which cannot be incident with both distinct
$b$-elements. It can contain at most one $c$-element. The empty flag
has four shadow elements, whereas any $c$-element has three. Neither
has the required shadow. Thus this residue fails the property.
\end{proof}


\section{Thin rank-three geometries and involution triples}
Throughout, the group preserves types and acts transitively on flags of
each type. Incidence is symmetric and reflexive, with no incident distinct
objects of one type. Every flag extends to a chamber. A geometry is thin
if every flag missing one type has exactly two extensions. Residual
connectedness requires the incidence graph of every residue of rank at
least two to be connected. Residual weak primitivity means that in every
nonempty residue at least one of its actions on a type is primitive.

\begin{lemma}\label{thin-triples}
Let a nonabelian finite simple group $G$ act as above on a thin, residually
connected geometry of rank three. For a chamber $C=\{p_0,p_1,p_2\}$,
write $H_i=G_{p_i}$ and $B=\bigcap_i H_i$. There are distinct involutions
$t_0,t_1,t_2$ such that
\[
 B=1,\qquad G=\langle t_0,t_1,t_2\rangle,\qquad
 H_i=\langle t_j:j\ne i\rangle,
 \qquad H_i\cap H_j=\langle t_k\rangle\quad(\{i,j,k\}=\{0,1,2\}).
\]
\end{lemma}
\begin{proof}
Put $R_i=\bigcap_{j\ne i}H_j$. Flag transitivity on the two extensions
of $C\setminus\{p_i\}$ gives $[R_i:B]=2$, so $B$ is normal in each $R_i$.

The chamber graph is connected. To see this, a rank-two residue has a
connected incidence graph, whose edges are its chambers; successive edges
sharing a vertex give a gallery. In rank three take a path of incident
objects between a vertex of one chamber and a vertex of another. Extend
each edge to a chamber. Consecutive chosen chambers contain a common
object, and the rank-two residue of that object supplies a gallery between
them. The end chambers connect in the same way. Thus all chambers are
connected by changes of one object.

The chambers adjacent to $C$ across its $i$-panel are images of $C$ under
$R_i$. Transporting along a gallery proves
$G=\langle R_0,R_1,R_2\rangle$: the generated subgroup reaches every
chamber, and it contains their common base stabilizer $B$. Therefore
$B\lhd G$. Since $B<R_i$, simplicity gives $B=1$. Each $R_i$ now has a
unique nonidentity element $t_i$, an involution. Within the residue of
$p_i$, the same gallery argument gives
$H_i=\langle R_j:j\ne i\rangle$. Distinctness follows because an element
in two $R_i$ lies in $B$. The stated intersections follow from the
stabilizer definitions.
\end{proof}

\begin{lemma}\label{thin-recognize}
Conversely take three distinct involutions generating $G$, and put
$H_i=\langle t_j:j\ne i\rangle$. Suppose the three intersections in
Lemma~\ref{thin-triples} hold. Their coset incidence system is a flag-transitive
geometry if and only if, for each ordering $i,j,k$ of the types,
\begin{equation}\label{products}
        (H_iH_j)\cap(H_iH_k)=H_i(H_j\cap H_k).
\end{equation}
It then is thin and residually connected. It is residually weakly primitive
if and only if at least one $H_i$ is maximal in $G$ and all three numbers
$|H_i|/2$ are prime.
\end{lemma}
\begin{proof}
Any incident pair of cosets has a common element, so $G$ is transitive
on incident pairs. Normalize the cosets of types $j,k$ to $H_j,H_k$.
The cosets $H_i g$ incident with both correspond exactly to the elements
of $(H_iH_j)\cap(H_iH_k)$. They extend to chambers in the base chamber's
orbit precisely when they meet $H_j\cap H_k$. This gives
\eqref{products}. Every flag then extends to such a chamber. The right
side is contained in the left side; since the triple intersection is
trivial, its size is $2|H_i|$. Thus equality is checked without constructing
an entire incidence graph.

The base chamber stabilizer is trivial and its panel stabilizers have
order two, so the geometry is thin. Each $H_i$ is generated by the two
panel stabilizers within it; the whole group is generated by all three.
Their coset incidence graphs are connected, proving residual connectedness.
Each rank-two residue is the polygon associated to a dihedral group of
order $2m_i$, where $m_i=|H_i|/2$. Its vertex stabilizers are reflections.
A reflection subgroup is maximal in a dihedral group of order $2m_i$
exactly when $m_i$ is prime: intermediate subgroups correspond to proper
divisors of the rotation order. Rank-one residues have degree two and
are primitive. The remaining condition is primitivity on some type of
the whole geometry, equivalent to maximality of its stabilizer $H_i$.
\end{proof}


\begin{thebibliography}{9}
\bibitem[JL]{jl} P.~Jacobs and D.~Leemans, An algorithmic analysis of the
intersection property, \emph{LMS J. Comput. Math.} 7 (2004), 284--299.
\bibitem[Cover]{cover} The project's accepted phase-four covering proof
and later four-lemma formal supplement. The latter proves only the
flag-compatible covering equivalence, not the coset computations.
\bibitem[D]{d} The project's previously accepted common finite-geometry
proofs and exact source-specific census packets. Each application names
which census completeness and group models it retains.
\bibitem[CC02]{cc} P.~J.~Cameron and P.~Cara, Independent generating sets
and geometries for symmetric groups, \emph{J. Algebra} 258 (2002), 641--650.
The transposition-tree model is established prior mathematics.
\end{thebibliography}
\end{document}
